% ----------------------------------------------------------
\chapter{Análise Comparativa e Decisão}
% ----------------------------------------------------------


Após a definição de dois possíveis traçados, a próxima etapa é a análise qualitativa 
e quantitativa de suas características para a escolha do mais adequado para se desenvolver 
o projeto rodoviário.  


\section{Análise Qualitativa}

Qualitativamente, os traçados foram escolhidos considerando a interação entre os perfis horizontal e vertical, conforme preconizado por \cite{lee2003}. Dessa forma, buscou-se compatibilizar ao máximo as curvas horizontais com as curvas verticais.

Observa-se que o Traçado B priorizou o menor corte de área verde, favorecendo a preservação ambiental. Em contrapartida, este traçado intercepta um maior número de zonas de proteção permanente em comparação ao Traçado A. Quanto à compatibilização das curvas verticais e horizontais, o Traçado B apresenta melhor desempenho neste quesito.



\section{Análise Qualitativa}

Já para a análise quantitativa, foram utilizados os critérios de comprimento 
total do traçado, movimentação de terra, corte de vegetação e comprimento do 
tabuleiro da ponte sobre o córrego, em que cada critério tem um peso específico 
que representa sua relevância na definição do traçado final (Tabela \ref{tab:tab_2}). Em mãos 
desses dados, preencheu-se a tabela abaixo: 



\begin{table}[h]
	
	\caption{Análise auxiliar dos traçados propostos}\label{tab:tab_2}
	\begin{tabular}{|ll|l|c|c|}
		\hline
		\multicolumn{2}{|l|}{Elementos}                                               & \textbf{Peso} & \multicolumn{1}{l|}{\textbf{Traçado A}} & \multicolumn{1}{l|}{\textbf{Traçado B}} \\ \hline
		\multicolumn{2}{|l|}{\textbf{1 -   Comprimento total do traçado (m)}}         & 0.15          & 5701.482                                & 5018.937                                \\ \hline
		\multicolumn{2}{|l|}{\textbf{2 -   Movimentação de Terras (m²)}}              & 0.45          & 273388.66                               & 153521                                  \\ \hline
		\multicolumn{2}{|l|}{\textbf{3 - Área da   obra de arte especial (m²)}}       & 0.15          & 43522.322                               & 1942.504                                \\ \hline
		\multicolumn{2}{|l|}{\textbf{4 - Área   referente ao impacto ambiental (m²)}} & 0.25          & 39594.691                               & 20926.51                                \\ \hline
		\multicolumn{2}{|l|}{\textbf{Soma}}                                           & \textbf{1}    & \textbf{1.4E+05}                        & \textbf{7.5E+04}                        \\ \hline
	\end{tabular}
\end{table}


\section{Escolha do traçado }


Pode-se perceber que o corte de vegetação e a movimentação de terra 
possuem peso muito maior que dos outros dois critérios, coerente com o impacto 
que esses fatores geram economicamente e ambientalmente na obra. O Traçado 
B possui uma movimentação de terra consideravelmente menor que a 
movimentação do outro traçado, que por sua vez possui menor corte de vegetação. 
Já o comprimento total e passagem pelo córrego são muito semelhantes.  
Ambos traçados estão de acordo com as recomendações para traçado em 
planta e perfil do DNER, portando a definição de qual alinhamento utilizar foi feito 
por base nos resultados da Tabela 4. A movimentação de terra foi o critério 
determinante para a decisão do traçado, justificando, portanto, a escolha do 
Traçado B como a melhor opção, já que ele apresenta significantemente 
menor volume de aterro. 





